\documentclass[11pt,a4paper]{article}

\usepackage[utf8]{inputenc}
\usepackage[margin=1in]{geometry}
\usepackage{hyperref}
\usepackage{booktabs}
\usepackage{xcolor}
\usepackage{listings}
\usepackage{tcolorbox}
\usepackage{enumitem}
\usepackage{titlesec}
\usepackage{fancyhdr}
\usepackage{microtype}
\usepackage{amsmath}

% Color Definitions
\definecolor{primaryblue}{RGB}{14, 90, 160}
\definecolor{accentorange}{RGB}{230, 95, 10}
\definecolor{darkgray}{RGB}{45, 52, 54}
\definecolor{lightgray}{RGB}{245, 246, 250}
\definecolor{codebg}{RGB}{240, 242, 245}
\definecolor{successgreen}{RGB}{39, 174, 96}

% Hyperlink Configuration
\hypersetup{
    colorlinks=true,
    linkcolor=primaryblue,
    filecolor=primaryblue,      
    urlcolor=primaryblue,
    citecolor=primaryblue
}

% Code Listing Configuration
\lstset{
    backgroundcolor=\color{codebg},
    basicstyle=\ttfamily\small\color{darkgray},
    keywordstyle=\bfseries\color{primaryblue},
    commentstyle=\itshape\color{gray},
    stringstyle=\color{accentorange},
    frame=single,
    framesep=4pt,
    rulecolor=\color{lightgray},
    breaklines=true,
    showstringspaces=false,
    tabsize=2
}

% Header and Footer
\pagestyle{fancy}
\fancyhf{}
\rhead{\color{gray}\small ARSpace Setup \& Deployment Guide}
\lhead{\color{gray}\small Commercial Real Estate AR Platform}
\cfoot{\thepage}

% Section formatting
\titleformat{\section}
  {\normalfont\Large\bfseries\color{primaryblue}}{\thesection}{1em}{}
\titleformat{\subsection}
  {\normalfont\large\bfseries\color{darkgray}}{\thesubsection}{1em}{}
\titleformat{\subsubsection}
  {\normalfont\normalsize\bfseries\color{darkgray}}{\thesubsubsection}{1em}{}

\begin{document}

% Title Block
\begin{titlepage}
    \centering
    \vspace*{1.5cm}
    
    {\Huge\bfseries\color{primaryblue} ARSpace: Augmented Reality Platform}\\[0.4cm]
    {\LARGE\bfseries Commercial Real Estate Layout Visualization}\\[1.5cm]
    
    \textbf{\Large Comprehensive Technical Manual}\\[0.3cm]
    {\large Unity 6 Setup, Project Configuration, Testing, and Android Deployment}\\[2cm]
    
    \begin{tcolorbox}[colback=lightgray,colframe=primaryblue,width=0.9\textwidth,arc=4mm]
        \centering
        \textbf{Document Version:} 1.0.0 \\
        \textbf{Target Platform:} Android (ARCore / ARM64 / OpenGLES3) \\
        \textbf{Engine Version:} Unity 6 (6000.5.6f1) / URP 17.5.0 / AR Foundation 6.5.0 \\
        \textbf{Repository Branch:} \texttt{feature/arspace-commercial-platform}
    \end{tcolorbox}
    
    \vfill
    {\large Prepared for Engineering, QA, and Deployment Teams}\\[0.5cm]
    {\large \today}
\end{titlepage}

\tableofcontents
\newpage

% -------------------------------------------------------------
\section{Executive Summary}
% -------------------------------------------------------------
\textbf{ARSpace} is an enterprise Augmented Reality (AR) application engineered for commercial real estate (CRE) brokers, workplace interior planners, and enterprise tenants. It facilitates real-time 3D spatial planning inside bare-shell corporate floor plates before construction begins.

The application incorporates an enterprise architectural pattern centered around:
\begin{itemize}[noitemsep]
    \item \textbf{Decoupled Event Bus (\texttt{GameEvents})}: Complete separation of user interface and AR subsystems.
    \item \textbf{Central Registry (\texttt{ServiceLocator})}: Lifecycle-managed services accessible across layers.
    \item \textbf{Spatial Anchor Clustering (\texttt{AnchorService})}: Groups items within 1.5\,m under native AR anchors to respect ARCore resource limits.
    \item \textbf{De-Duplicated Floor Area Calculation (\texttt{FloorAreaCalculator})}: 2D spatial grid sampling (0.25\,m resolution) that prevents overlap double-counting.
    \item \textbf{Origin-Relative Persistence (\texttt{LayoutOriginManager})}: Offline layout serialization restored at any physical floor location.
    \item \textbf{Executive Presentation Mode (\texttt{PresentationModeController})}: Immersion toggle for uncluttered client walkthroughs.
\end{itemize}

---

% -------------------------------------------------------------
\section{System Requirements and Prerequisites}
% -------------------------------------------------------------

\subsection{Workstation Development Environment}
\begin{itemize}
    \item \textbf{Operating System:} macOS 13+ (Apple Silicon or Intel) or Windows 11 (64-bit).
    \item \textbf{Unity Hub:} Latest version (3.8.0 or newer).
    \item \textbf{Unity Editor:} Version \textbf{6000.5.6f1} (Unity 6).
    \item \textbf{Unity Build Modules Required:}
    \begin{itemize}[noitemsep]
        \item \texttt{Android Build Support}
        \item \texttt{Android SDK \& NDK Tools}
        \item \texttt{OpenJDK}
    \end{itemize}
\end{itemize}

\subsection{Target Android Device Requirements}
\begin{itemize}
    \item \textbf{Hardware:} ARCore-certified Android device (e.g., Google Pixel 4+, Samsung Galaxy S10+, OnePlus 7+, etc.).
    \item \textbf{Operating System:} Android 8.0 (Oreo / API level 26) or higher.
    \item \textbf{CPU Architecture:} 64-bit ARM (\texttt{arm64-v8a}).
    \item \textbf{Google Play Services for AR:} Installed and updated to the latest version via the Google Play Store.
    \item \textbf{Hardware Connection:} USB-C to USB-A/C data cable supporting ADB debugging.
\end{itemize}

---

% -------------------------------------------------------------
\section{Repository Setup and Initial Project Load}
% -------------------------------------------------------------

\subsection{Cloning the Project}
Open terminal and clone the repository, switching to the verified feature branch:
\begin{lstlisting}[language=bash]
# Clone repository
git clone https://github.com/KodeWithKeshav/ARSpace.git
cd ARSpace

# Check out the complete, production-ready feature branch
git checkout feature/arspace-commercial-platform
\end{lstlisting}

\subsection{Opening in Unity Hub}
\begin{enumerate}
    \item Launch \textbf{Unity Hub}.
    \item Navigate to \textbf{Projects} $\rightarrow$ click \textbf{Add} $\rightarrow$ \textbf{Add project from disk}.
    \item Select the inner Unity project directory:
    \begin{lstlisting}[language=bash]
    ARSpace/UnityProject/ARSpace
    \end{lstlisting}
    \item Verify that the Editor Version is set to \textbf{Unity 6 (6000.5.6f1)}.
    \item Click on the project name to launch the Editor.
    \item Allow Unity several minutes on first launch to import assets, compile package caches, and index ScriptableObjects.
\end{enumerate}

---

% -------------------------------------------------------------
\section{Automated One-Click Project Configuration}
% -------------------------------------------------------------
To ensure zero manual configuration errors, ARSpace includes custom idempotent Editor builders exposed under the top menu bar: \textbf{\texttt{ARSpace}}.

\begin{figure}[h!]
\centering
\begin{tcolorbox}[colback=white,colframe=primaryblue,title=\textbf{The ARSpace Editor Menu Hierarchy},arc=2mm]
\begin{itemize}[leftmargin=1.5em]
    \item \textbf{\texttt{ARSpace $\rightarrow$ Validate Project}}: Comprehensive audit of all settings, assets, and shaders.
    \item \textbf{\texttt{ARSpace $\rightarrow$ Apply Android AR Build Settings}}: One-click enforcer for IL2CPP, ARM64, and OpenGLES3.
    \item \textbf{\texttt{ARSpace $\rightarrow$ Setup Plane Grid Assets}}: Generates materials and visualizer prefabs.
    \item \textbf{\texttt{ARSpace $\rightarrow$ Rebuild Furniture Catalogue}}: Scans 52 3D models and builds runtime assets.
\end{itemize}
\end{tcolorbox}
\end{figure}

Execute the setup steps in the exact sequence described below:

\subsection{Step 1: Enforce Android AR Build Settings}
Select \textbf{\texttt{ARSpace $\rightarrow$ Apply Android AR Build Settings}}.
This script executes \texttt{BuildSettingsEnforcer.cs} and automatically sets:
\begin{itemize}[noitemsep]
    \item Build Target: \texttt{Android}
    \item Scripting Backend: \texttt{IL2CPP}
    \item Target Architecture: \texttt{ARM64} (\texttt{AndroidArchitecture.ARM64})
    \item Graphics API: \texttt{OpenGLES3} (Vulkan is explicitly disabled to prevent known driver crashes in ARCore on mid-range Mali/Adreno GPUs)
    \item Colour Space: \texttt{Linear}
    \item Minimum API Level: \texttt{26} (Android 8.0 Oreo)
    \item Scene at Build Index 0: \texttt{Assets/Scenes/ARWorkspace.unity}
\end{itemize}

\subsection{Step 2: Generate Plane Grid Visualization Assets}
Select \textbf{\texttt{ARSpace $\rightarrow$ Setup Plane Grid Assets}}.
This generates:
\begin{itemize}[noitemsep]
    \item \texttt{Assets/Art/Materials/M\_PlaneGrid.mat} bound to \texttt{ARSpace/PlaneGrid.shader}.
    \item \texttt{Assets/Art/Materials/M\_PlaneOutline.mat} bound to \texttt{ARSpace/UnlitLine.shader}.
    \item Configures the \texttt{Plane.prefab} with \texttt{PlaneGridVisualizer} and \texttt{ARPlaneMeshVisualizer}.
\end{itemize}

\subsection{Step 3: Rebuild Furniture Catalogue}
Select \textbf{\texttt{ARSpace $\rightarrow$ Rebuild Furniture Catalogue}}.
This automated pipeline processes all 52 models in \texttt{Assets/Art/Models/}:
\begin{enumerate}[noitemsep]
    \item Normalizes 3D pivots to floor level ($Y = 0$).
    \item Attaches \texttt{PlacedObject} and tight-fitting \texttt{BoxCollider}.
    \item Renders 256$\times$256 offscreen studio preview thumbnails as sprites under \texttt{Assets/Art/Thumbnails/}.
    \item Creates \texttt{FurnitureItem} ScriptableObjects with measured bounds, real-world footprints, clearance margins, and seating capacities.
    \item Updates \texttt{Assets/ScriptableObjects/FurnitureDatabase.asset} with zero duplicate IDs.
\end{enumerate}

\subsection{Step 4: Execute Project Validator}
Select \textbf{\texttt{ARSpace $\rightarrow$ Validate Project}}.
A GUI window will display. Confirm that all 35+ verification checks display a green checkmark (\textcolor{successgreen}{\textbf{\checkmark}}):
\begin{itemize}[noitemsep]
    \item Build target, scripting backend, architecture, graphics API
    \item URP asset, XR loader, AR session profile
    \item Plane visualizer shaders and materials
    \item Catalogue items and database lookup tables
    \item Placement, selection, persistence, and space analytics scripts
    \item Automated test runner assembly
\end{itemize}

---

% -------------------------------------------------------------
\section{Running and Testing in Unity Editor}
% -------------------------------------------------------------

\subsection{Opening the Main Scene}
In the Project window, double-click:
\begin{lstlisting}
Assets/Scenes/ARWorkspace.unity
\end{lstlisting}
Verify that the scene hierarchy contains:
\begin{itemize}[noitemsep]
    \item \textbf{\texttt{ARSpaceApp}}: Core state machine coordinator.
    \item \textbf{\texttt{XR Origin (XR Rig)}}: Contains \texttt{ARCamera}, \texttt{ARPlaneManager}, \texttt{ARAnchorManager}, and \texttt{ARRaycastManager}.
    \item \textbf{\texttt{ServiceContainer}}: Holds \texttt{CatalogService}, \texttt{ARPlacementManager}, \texttt{ObjectSelectionService}, \texttt{LayoutOriginManager}, and \texttt{SpaceAnalyticsService}.
    \item \textbf{\texttt{UI\_Canvas}}: Contains the Onboarding Coach, Reticle, Selection Toolbar, Layout Menu Panel, and Analytics Panel.
\end{itemize}

\subsection{Running the Automated Test Suite}
ARSpace includes a comprehensive EditMode test suite.
\begin{enumerate}
    \item Open the Unity Test Runner via \textbf{Window $\rightarrow$ General $\rightarrow$ Test Runner}.
    \item Select the \textbf{EditMode} tab.
    \item Click \textbf{Run All}.
    \item Verify that all test fixtures pass:
    \begin{itemize}[noitemsep]
        \item \texttt{FurnitureDatabaseTests}: $O(1)$ dictionary lookups, category filtering, case-insensitivity.
        \item \texttt{LayoutPersistenceTests}: JSON serialization, preset loading, missing ID tolerance.
        \item \texttt{SpaceAnalyticsTests}: Circulation ratios, seating density thresholds, zero-division safety.
        \item \texttt{ServiceLocatorTests}: Register, retrieve, and lifecycle cleanup.
        \item \texttt{AnchorClusteringTests}: 1.5\,m spatial cluster evaluation and 20-anchor cap enforcement.
    \end{itemize}
\end{enumerate}

\subsection{Simulating in Editor Play Mode}
Click the \textbf{Play} button in Unity Editor.
\begin{itemize}
    \item Unity's XR Simulation subsystem will load a simulated room environment.
    \item Move the camera using \texttt{W/A/S/D} keys and right-mouse button look.
    \item The simulated floor plane will be detected, showing the orange anti-aliased grid and radar pulse.
    \item Select an item from the catalogue drawer and left-click on the floor to test placement and manipulation.
\end{itemize}

---

% -------------------------------------------------------------
\section{Android Device Configuration}
% -------------------------------------------------------------

\subsection{Enabling Developer Options on Android}
\begin{enumerate}
    \item On your Android device, open \textbf{Settings $\rightarrow$ About phone}.
    \item Scroll down to \textbf{Build number}.
    \item Tap \textbf{Build number} rapidly \textbf{7 times} until you see the prompt: \textit{"You are now a developer!"}
    \item Return to \textbf{Settings $\rightarrow$ System $\rightarrow$ Developer options}.
    \item Enable the toggle for \textbf{Developer options}.
    \item Scroll down and enable \textbf{USB debugging}.
\end{enumerate}

\subsection{Installing Google Play Services for AR (ARCore)}
\begin{enumerate}
    \item Open the \textbf{Google Play Store} on your Android device.
    \item Search for \textbf{Google Play Services for AR}.
    \item If not installed, tap \textbf{Install}. If already installed, tap \textbf{Update} to ensure compatibility with AR Foundation 6.5.
\end{enumerate}

\subsection{Verifying ADB Communication}
Connect the device to your computer using a USB-C data cable. On your device, accept the prompt: \textit{"Allow USB debugging from this computer?"} (check \textit{"Always allow"}).

Open a terminal on your computer and run:
\begin{lstlisting}[language=bash]
adb devices
\end{lstlisting}
Confirm that your device is listed with state \texttt{device}:
\begin{lstlisting}
List of devices attached
1234567890ABCDEF    device
\end{lstlisting}

---

% -------------------------------------------------------------
\section{Building and Deploying to Android}
% -------------------------------------------------------------

\subsection{Method 1: Direct "Build and Run" (Recommended)}
This method compiles the IL2CPP binaries, generates the APK, deploys it to the connected phone, and automatically launches the app.

\begin{enumerate}
    \item Connect your Android device via USB and verify ADB recognition.
    \item In Unity Editor, open \textbf{File $\rightarrow$ Build Profiles} (or \textbf{Build Settings}).
    \item Select \textbf{Android} platform.
    \item In the \textbf{Run Device} dropdown, select your connected Android device model.
    \item Verify that \texttt{Assets/Scenes/ARWorkspace.unity} is checked at index 0.
    \item Click \textbf{Build and Run}.
    \item Choose a location to save the output file (e.g., \texttt{ARSpace/builds/ARSpace.apk}).
    \item Unity will compile the IL2CPP C++ code, pack the APK, install it over ADB, and launch the application on your phone.
\end{enumerate}

\subsection{Method 2: Standalone APK Export and Sideload}
If building for distribution to clients or other team members:

\begin{enumerate}
    \item In Unity, open \textbf{File $\rightarrow$ Build Settings}.
    \item Click \textbf{Build} (not Build and Run).
    \item Save the file as \texttt{ARSpace\_Release.apk}.
    \item Once the build completes, install it manually via command line:
    \begin{lstlisting}[language=bash]
    # Sideload the APK onto the connected Android device
    adb install -r builds/ARSpace_Release.apk
    \end{lstlisting}
    \item The app icon labeled \textbf{ARSpace} will appear on the device home screen.
\end{enumerate}

---

% -------------------------------------------------------------
\section{In-App Mobile Operating Guide}
% -------------------------------------------------------------

When the application launches on your Android device, follow this walkthrough workflow:

\subsection{1. Floor Mapping \& Room Onboarding}
\begin{itemize}
    \item When prompted, grant camera permissions.
    \item The \textbf{Onboarding Coach} banner at the top reads: \textit{"Map Your Workspace: Slowly sweep your phone over the bare floor to detect surfaces."}
    \item Pan the device camera smoothly across the room. Horizontal floor surfaces will light up with an anti-aliased 0.5\,m primary / 1.0\,m secondary grid and a subtle radar sweep.
    \item Once $\ge 1.5\,\text{m}^2$ of floor is detected, the coach automatically advances to step 2.
\end{itemize}

\subsection{2. Catalogue Selection and Placement}
\begin{itemize}
    \item The bottom catalogue drawer categorizes items across 8 functional areas: Workstations, Cabins, Conference, Cafeteria, Reception, Partitions, Decor, and Equipment.
    \item Tap any item thumbnail (e.g., \textbf{6-Person Team Pod Cubicle}).
    \item Aim your device at the floor. The reticle will project the real-world rectangular footprint of the furniture item on the floor.
    \item Tap the primary \textbf{Place} button (or tap directly on the reticle).
    \item The furniture item instantiates at full 1:1 scale, automatically snapped to face the user, and parents to a clustered native AR anchor.
\end{itemize}

\subsection{3. Object Selection and Touch Manipulation}
\begin{itemize}
    \item \textbf{Selection:} Tap any placed furniture piece. An active boundary ring highlights the item's perimeter and its required circulation clearance margin.
    \item \textbf{Translation (1-Finger Drag):} Drag the object across the floor. Translation is strictly constrained to the detected floor surface.
    \item \textbf{Collision Feedback:} If you drag an object into another object's clearance boundary, both items dynamically tint orange, alerting you to circulation conflicts.
    \item \textbf{Rotation (2-Finger Twist):} Twist with two fingers to rotate the item in place. Rotation features $5^\circ$ magnetic snapping for effortless alignment with office walls.
    \item \textbf{Contextual Toolbar:} Floating above the selected item is a toolbar providing:
    \begin{itemize}[noitemsep]
        \item \textbf{Duplicate}: Spawns an identical item with identical rotation offset by 0.5\,m.
        \item \textbf{Rotate 90$^\circ$}: Quick orthogonal orientation.
        \item \textbf{Lock/Unlock}: Locks translation to prevent accidental touches during reviews.
        \item \textbf{Delete}: Removes the item and cleanly de-allocates its native anchor.
        \item \textbf{CRE Specs}: Displays item name, dimensions ($L \times W \times H$), and seat count.
    \end{itemize}
\end{itemize}

\subsection{4. Commercial Space Analytics Dashboard}
\begin{itemize}
    \item Tap the \textbf{Analytics} tab to expand the live CRE dashboard bottom sheet.
    \item \textbf{Usable Floor Area}: De-duplicated room square meterage calculated via 2D spatial occupancy grid sampling.
    \item \textbf{Footprint vs. Circulation Ratio}: Live percentage of floor occupied by furniture vs. free egress space.
    \item \textbf{Compliance Badges}:
    \begin{itemize}[noitemsep]
        \item \textcolor{successgreen}{\textbf{COMPLIANT ($\ge 45\%$)}}: Standard corporate accessibility and egress compliance.
        \item \textcolor{accentorange}{\textbf{CONSTRAINED (30--45\%)}}: High-density team pod zoning.
        \item \textcolor{red}{\textbf{NON-COMPLIANT ($< 30\%$)}}: Fire egress hazard alert.
    \end{itemize}
    \item \textbf{Seating Density}: Seats per 100\,m$^2$ with health categorization (Sparse, Optimal, High-Density, Overcrowded).
    \item \textbf{Unit Toggle}: Tap the unit toggle button to instantly convert between \textbf{Metric (m$^2$)} and \textbf{Imperial (sq ft)}.
\end{itemize}

\subsection{5. Layout Persistence and Built-in Presets}
\begin{itemize}
    \item Open the \textbf{Layout Menu}.
    \item \textbf{Saving Custom Layouts:} Enter a name (e.g., \texttt{"Tech\_Floor\_Plan"}) and tap \textbf{Save}. Layout coordinates are atomically written to local disk relative to the room's physical origin pose.
    \item \textbf{Restoring Layouts:} Tap any saved layout or choose from the 3 built-in commercial templates:
    \begin{enumerate}[noitemsep]
        \item \textbf{Open Plan Team Pod (24 Seats)}: 4x 6-person cubicles, printers, water units, whiteboards.
        \item \textbf{Hybrid Team Zone (16 Seats)}: Collaborative cubicles, phone booths, huddle desk, sofa.
        \item \textbf{Executive Floor \& Boardroom (22 Seats)}: CEO suite, cabin, 12-person boardroom table, reception.
    \end{enumerate}
    \item The entire layout reconstructs instantly at your verified floor location.
\end{itemize}

\subsection{6. Executive Presentation Mode}
\begin{itemize}
    \item During formal client walkthroughs, tap the \textbf{Present} button.
    \item All UI chrome, drawers, reticle, and plane grid lines are hidden, leaving an uninterrupted, photorealistic view of the office space.
    \item Soft contact ground shadows render under furniture, anchored to real floor surfaces via \texttt{ShadowReceiver.shader}.
    \item Real-time ambient lighting and directional light intensity adapt smoothly via \texttt{LightEstimationBinder}.
    \item To exit, tap the subtle floating restore icon in the corner or \textbf{double-tap anywhere on the screen}.
\end{itemize}

---

% -------------------------------------------------------------
\section{Troubleshooting and Common Solutions}
% -------------------------------------------------------------

\begin{table}[h!]
\centering
\small
\begin{tabular}{p{4.5cm} p{4.5cm} p{5.5cm}}
\toprule
\textbf{Symptom} & \textbf{Probable Cause} & \textbf{Resolution} \\
\midrule
App crashes immediately on launch with black screen. & Vulkan graphics API enabled or ARCore missing. & Run \texttt{ARSpace $\rightarrow$ Apply Android AR Build Settings} to enforce OpenGLES3. Ensure \textit{Google Play Services for AR} is updated. \\
\addlinespace
Floor planes take a long time to detect. & Low texture / glossy floor surface or poor lighting. & Pan device slowly in a figure-8 motion. Ensure room has adequate lighting and visible floor texture. \\
\addlinespace
\texttt{adb devices} shows \texttt{unauthorized}. & USB debugging RSA key not yet accepted on device. & Re-plug USB cable and tap \textit{"Always allow from this computer"} on the phone. \\
\addlinespace
Build fails with NDK or IL2CPP toolchain error. & Incomplete Android SDK/NDK installation. & Open Unity Hub $\rightarrow$ Installs $\rightarrow$ Unity 6 $\rightarrow$ Add Modules $\rightarrow$ check \textit{Android SDK \& NDK Tools} and \textit{OpenJDK}. \\
\addlinespace
Catalogue shows missing models or pink thumbnails. & Catalogue database has not been compiled. & In Unity Editor, execute \texttt{ARSpace $\rightarrow$ Rebuild Furniture Catalogue}. \\
\bottomrule
\end{tabular}
\end{table}

\vfill
\begin{center}
\color{gray}\small --- End of ARSpace Technical Manual ---
\end{center}

\end{document}
